%% kapa: exam=1 round=25 year=2014 period=2 form=A subject=appraisal_law subject_ko=부동산관계법규 seq=17 number=17
%% source: Q-Net 감정평가사 기출문제 (cst00309) artlSeq=5101207 file=extracted/1008/5101207/2014년 제25회 감정평가사 1차시험 기출문제/2교시(A형).hwp
%% conversion: kordoc 4.24.51 -> tools/gen1.py (auto); review=proofread-sonnet; flags=none
\begin{question}{17}
부동산 가격공시 및 감정평가에 관한 법령상 지가의 공시에 관한 설명으로 ( )에 알맞은 것은?
\begin{framedtext}
ㆍ정부는 지가공시의 주요사항에 관한 보고서를 매년 정기국회의 ( ㄱ ) 전까지 국회에 제출하여야 한다. \par
ㆍ시장이 개별공시지가확인서를 발급하는 때에는 당해 시의 ( ㄴ )(으)로 정하는 바에 따라 신청인으로부터 소정의 수수료를 받을 수 있다. \par
\end{framedtext}
\begin{choices}
  \item \makebox[9em][l]{ㄱ: 개회 30일}ㄴ: 규칙
  \item \makebox[9em][l]{ㄱ: 개회}ㄴ: 규칙
  \item \makebox[9em][l]{ㄱ: 개회 30일}ㄴ: 조례
  \item \makebox[9em][l]{ㄱ: 개회}ㄴ: 조례
  \item \makebox[9em][l]{ㄱ: 폐회}ㄴ: 조례
\end{choices}
\end{question}
\answer{4}
